\documentclass[12pt,a4paper]{ctexart}

\usepackage[margin=1in]{geometry}
\usepackage{array}
\usepackage{booktabs}
\usepackage{hyperref}
\usepackage{longtable}
\usepackage{tabularx}
\usepackage{xcolor}
\usepackage{listings}
\usepackage{float}

\hypersetup{
  colorlinks=true,
  linkcolor=blue,
  urlcolor=blue,
  citecolor=blue
}

\lstdefinestyle{fighterstyle}{
  basicstyle=\ttfamily\small,
  breaklines=true,
  columns=fullflexible,
  frame=single,
  backgroundcolor=\color{gray!5}
}

\title{Final Project Design Document}
\author{}
\date{2026-04-17}

\begin{document}

\maketitle

本文档基于仓库当前实现整理，内容参考提供的 \texttt{Design Doc.docx} 框架，但所有描述都以当前代码分支的真实实现状态为准。文中会明确区分“已实现功能”“软件侧已冻结接口”和“后续目标架构”，避免把尚未落地的设计目标写成现状。

\section{引言与概述}

\subsection{项目摘要}

本项目是在 \texttt{DE1-SoC} 平台上实现一个双人格斗游戏原型。当前仓库已经完成一条可运行的 \texttt{Phase 1} 主链路：

\begin{itemize}
  \item \texttt{HPS} 侧 C 程序负责双人输入解析、游戏状态机、碰撞与伤害计算、音频事件调度。
  \item 显示目前主要通过 \texttt{HPS} 侧软件渲染完成，优先写 \texttt{/dev/fb0}，失败时回退到控制台文本输出。
  \item \texttt{FPGA} 侧已经实现 \texttt{WM8731} 音频外设 \texttt{fighter\_audio\_wm8731}，通过 lightweight HPS-to-FPGA bridge 暴露 \texttt{Avalon-MM} 寄存器。
  \item 代码中已经预留一套 \texttt{fighter\_mmio} 游戏状态寄存器协议，供未来的 \texttt{FPGA VGA} 渲染外设接入，但当前分支还没有对应的 VGA 自定义渲染模块。
\end{itemize}

\subsection{当前实现与目标架构的关系}

\begin{longtable}{>{\raggedright\arraybackslash}p{0.28\textwidth} >{\raggedright\arraybackslash}p{0.18\textwidth} >{\raggedright\arraybackslash}p{0.42\textwidth}}
\toprule
层级 & 当前状态 & 说明 \\
\midrule
\endhead
双 USB 键盘输入 & 已实现 & \texttt{libusb} 轮询最多两个 HID boot keyboard \\
战斗逻辑状态机 & 已实现 & 菜单、对局、结算、攻击判定、格挡、超时、KO \\
视频输出 & 已实现，但在 HPS 侧 & 软件 framebuffer/console 渲染 \\
FPGA 音频外设 & 已实现 & \texttt{fighter\_audio\_wm8731} + \texttt{WM8731} bring-up \\
游戏状态 MMIO 编码 & 已实现 & \texttt{fighter\_mmio\_encode()} 可输出 22 个 32-bit 寄存器 \\
FPGA VGA 格斗渲染外设 & 未实现 & 模板中的理想化硬件路径目前尚未落地 \\
Linux kernel driver & 未实现 & 当前直接使用 \texttt{/dev/fb0} 与 \texttt{/dev/mem} \\
\bottomrule
\end{longtable}

\subsection{设计假设与约束}

\begin{itemize}
  \item 目标板卡为 \texttt{DE1-SoC}。
  \item 默认最多接入 2 个 USB 键盘。
  \item 默认输入映射为 \texttt{W/A/S/D/J/K/L}。
  \item 主逻辑帧率按 60 FPS 运行。
  \item 当前 \texttt{Quartus} 顶层是 \texttt{audio-only build}，\texttt{VGA} 顶层管脚在 \texttt{hw/soc\_system\_top.sv} 中被静态拉低。
\end{itemize}

\section{顶层系统框图}

\subsection{当前真实数据路径}

\begin{lstlisting}[style=fighterstyle]
USB Keyboard 1/2
      |
      v
usb_hid_keyboard.c  (libusb, up to 2 devices)
      |
      v
fighter_input.c
      |
      v
fighter_game.c  ------------------------------+
      |                                       |
      | game state                            | audio events
      v                                       v
fighter_renderer.c                    fighter_audio.c
      |                                       |
      |                                       +--> command backend fallback
      |                                            (aplay / ffplay / afplay)
      v
/dev/fb0 or console                         /dev/mem + lw bridge
                                                  |
                                                  v
                                       fighter_audio_wm8731 (FPGA)
                                                  |
                                                  v
                                               WM8731
\end{lstlisting}

\subsection{目标扩展路径}

模板中的目标架构仍然有价值，但在当前项目中应理解为下一阶段扩展路径：

\begin{lstlisting}[style=fighterstyle]
fighter_game.c
    |
    v
fighter_mmio_encode()
    |
    v
HPS-to-FPGA lightweight bridge
    |
    v
Future fighter VGA peripheral
    |
    v
VGA timing / sprite / UI output
\end{lstlisting}

\subsection{关键接口说明}

\begin{itemize}
  \item USB 输入接口：\texttt{usb\_hid\_keyboard\_manager\_init()} 会枚举 \texttt{HID boot keyboard}，寻找 interrupt IN endpoint，并轮询最多两个设备。
  \item 视频接口：当前不是自定义 FPGA VGA 协议，而是 HPS 用户态直接访问 \texttt{/dev/fb0}。因此模板中提到的 \texttt{HSYNC/VSYNC/RGB} 渲染协议，当前分支只存在于板级顶层端口，不存在项目自定义 VGA 模块实现。
  \item 音频接口：当前真实硬件接口为 \texttt{Avalon-MM + WM8731 serial audio + I2C}。HPS 通过物理地址 \texttt{0xFF203040} 访问音频外设，FPGA 负责 FIFO、I2C 初始化和串行发送。
\end{itemize}

\section{自定义外设寄存器映射}

本项目目前有两套寄存器协议：

\begin{enumerate}
  \item 已经在 FPGA 上实现并可上板验证的 \texttt{audio MMIO} 协议。
  \item 已经在软件侧固定、但尚未连接到 FPGA 渲染硬件的 \texttt{fighter MMIO} 协议。
\end{enumerate}

\subsection{已实现的音频 MMIO 协议}

默认物理地址为 \texttt{0xFF203040}，由 \texttt{sw/audio/fighter\_audio.c} 与 \texttt{hw/fighter\_audio.sv} 共同约定。

\begin{longtable}{>{\raggedright\arraybackslash}p{0.12\textwidth} >{\raggedright\arraybackslash}p{0.18\textwidth} >{\raggedright\arraybackslash}p{0.08\textwidth} >{\raggedright\arraybackslash}p{0.48\textwidth}}
\toprule
Offset & Name & R/W & 说明 \\
\midrule
\endhead
\texttt{0x00} & \texttt{control} & RW & 写 \texttt{bit2} 清读 FIFO，写 \texttt{bit3} 清写 FIFO；读回状态字 \\
\texttt{0x04} & \texttt{fifospace} & R & \texttt{[31:24]} 左声道可写空间，\texttt{[23:16]} 右声道可写空间 \\
\texttt{0x08} & \texttt{leftdata} & W & 左声道采样，软件写入 \texttt{int16 << 16} \\
\texttt{0x0C} & \texttt{rightdata} & W & 右声道采样，软件写入 \texttt{int16 << 16} \\
\bottomrule
\end{longtable}

\noindent \texttt{control} 读状态字的当前实现含义如下：

\begin{itemize}
  \item \texttt{bit0}: \texttt{codec\_init\_done}
  \item \texttt{bit1}: \texttt{codec\_init\_error}
  \item \texttt{bit2}: \texttt{tx\_underflow\_seen}
  \item \texttt{bit3}: \texttt{write\_overflow\_seen}
  \item \texttt{bits[15:8]}: \texttt{right\_count}
  \item \texttt{bits[23:16]}: \texttt{left\_count}
\end{itemize}

\subsection{软件侧已冻结的 fighter MMIO 协议}

这套协议由 \texttt{sw/render\_if/fighter\_mmio.c} 生成，并由 \texttt{sw/tests/test\_phase1.c} 校验，但当前分支还没有与之匹配的 FPGA VGA 外设。

总长度为 22 个 32-bit 字，地址空间共 \texttt{0x58} 字节。

\begin{longtable}{>{\raggedright\arraybackslash}p{0.12\textwidth} >{\raggedright\arraybackslash}p{0.24\textwidth} >{\raggedright\arraybackslash}p{0.16\textwidth} >{\raggedright\arraybackslash}p{0.34\textwidth}}
\toprule
Offset & 寄存器 & 类型 & 含义 \\
\midrule
\endhead
\texttt{0x00} & \texttt{GAME\_STATE} & Contracted & \texttt{bit[1:0]} 为 \texttt{MENU/PLAYING/GAME\_OVER}，\texttt{bit8} 为菜单动画帧，\texttt{bit9} 为 game over 是否可重新开始 \\
\texttt{0x04} & \texttt{PLAYER1\_X} & Contracted & 玩家 1 左上角 \texttt{x} \\
\texttt{0x08} & \texttt{PLAYER1\_Y} & Contracted & 玩家 1 左上角 \texttt{y} \\
\texttt{0x0C} & \texttt{PLAYER1\_STATE} & Contracted & \texttt{fighter\_visual\_state\_t} \\
\texttt{0x10} & \texttt{PLAYER2\_X} & Contracted & 玩家 2 左上角 \texttt{x} \\
\texttt{0x14} & \texttt{PLAYER2\_Y} & Contracted & 玩家 2 左上角 \texttt{y} \\
\texttt{0x18} & \texttt{PLAYER2\_STATE} & Contracted & \texttt{fighter\_visual\_state\_t} \\
\texttt{0x1C} & \texttt{PLAYER1\_HP} & Contracted & 玩家 1 剩余血量 \\
\texttt{0x20} & \texttt{PLAYER2\_HP} & Contracted & 玩家 2 剩余血量 \\
\texttt{0x24} & \texttt{ROUND\_TIMER} & Contracted & 剩余秒数，由 \texttt{fighter\_game\_round\_seconds\_remaining()} 生成 \\
\texttt{0x28} & \texttt{WINNER} & Contracted & \texttt{fighter\_winner\_t} \\
\texttt{0x2C} & \texttt{PLAYER1\_FACING} & Contracted & \texttt{1=朝右}，\texttt{0=朝左} \\
\texttt{0x30} & \texttt{PLAYER2\_FACING} & Contracted & \texttt{1=朝右}，\texttt{0=朝左} \\
\texttt{0x34} & \texttt{DEBUG\_FLAGS} & Contracted & \texttt{[7:0]} P1 最后攻击，\texttt{[15:8]} P2 最后攻击，\texttt{[19:16]} P1 attack phase，\texttt{[23:20]} P2 attack phase \\
\texttt{0x38} & \texttt{PLAYER1\_ATTACK\_CMD} & Contracted & \texttt{fighter\_attack\_command\_t} \\
\texttt{0x3C} & \texttt{PLAYER2\_ATTACK\_CMD} & Contracted & \texttt{fighter\_attack\_command\_t} \\
\texttt{0x40} & \texttt{PLAYER1\_STATE\_FRAME} & Contracted & 当前可视状态持续帧数 \\
\texttt{0x44} & \texttt{PLAYER2\_STATE\_FRAME} & Contracted & 当前可视状态持续帧数 \\
\texttt{0x48} & \texttt{PLAYER1\_EVENT\_FLAGS} & Contracted & 一帧脉冲事件，如 \texttt{ATTACK\_START/HIT/BLOCK/LAND/KO} \\
\texttt{0x4C} & \texttt{PLAYER2\_EVENT\_FLAGS} & Contracted & 同上 \\
\texttt{0x50} & \texttt{PLAYER1\_COMBAT\_RESULT} & Contracted & \texttt{NONE/HIT/BLOCKED/TRADE/WHIFF} \\
\texttt{0x54} & \texttt{PLAYER2\_COMBAT\_RESULT} & Contracted & 同上 \\
\bottomrule
\end{longtable}

\section{详细硬件设计}

\subsection{顶层硬件结构}

项目顶层为 \texttt{hw/soc\_system\_top.sv}，当前与项目直接相关的连线有：

\begin{itemize}
  \item \texttt{soc\_system} 实例负责 \texttt{HPS DDR3}、\texttt{USB}、\texttt{I2C}、\texttt{UART} 等 SoC 骨架。
  \item \texttt{audio\_*} 接口从 \texttt{soc\_system} 导出到板级 \texttt{AUD\_*} 管脚。
  \item \texttt{audio\_init\_done} 和 \texttt{audio\_init\_error} 映射到 \texttt{LEDR[0]} 和 \texttt{LEDR[1]}。
  \item \texttt{VGA} 管脚在当前构建中全部静态输出 0，表明该分支尚未接入项目定制的视频硬件。
\end{itemize}

\subsection{\texttt{fighter\_audio\_wm8731} 内部框图}

\begin{lstlisting}[style=fighterstyle]
Avalon-MM Slave
  |
  +--> control / status
  +--> left FIFO
  +--> right FIFO
          |
          v
     sample serializer ------> AUD_DACDAT
          |
          +--> clock dividers ---> AUD_XCK / AUD_BCLK / AUD_DACLRCK / AUD_ADCLRCK

I2C init FSM -----------------> FPGA_I2C_SCLK / FPGA_I2C_SDAT
                                |
                                v
                              WM8731
\end{lstlisting}

\subsection{音频外设的时钟与协议}

\texttt{fighter\_audio\_wm8731} 在单一 50 MHz 时钟域工作，注释中给出的默认派生时钟为：

\begin{itemize}
  \item \texttt{aud\_xck = clk / 4 = 12.5 MHz}
  \item \texttt{aud\_bclk = clk / 16 = 3.125 MHz}
  \item \texttt{lrck = clk / 1024 \approx 48.828 kHz}
\end{itemize}

Codec 初始化配置为：

\begin{itemize}
  \item 左对齐 \texttt{left-justified}
  \item 16-bit samples
  \item slave mode
  \item DAC playback enabled
\end{itemize}

\subsection{Platform Designer 组件接口}

\texttt{hw/fighter\_audio\_hw.tcl} 已经把该模块包装成 Qsys/Platform Designer 组件，暴露以下接口：

\begin{itemize}
  \item \texttt{clock}
  \item \texttt{reset}
  \item \texttt{avalon\_slave\_0}
  \item \texttt{audio} conduit
\end{itemize}

其中 \texttt{avalon\_slave\_0} 端口包括：

\begin{itemize}
  \item \texttt{avs\_chipselect}
  \item \texttt{avs\_read}
  \item \texttt{avs\_write}
  \item \texttt{avs\_address[1:0]}
  \item \texttt{avs\_writedata[31:0]}
  \item \texttt{avs\_readdata[31:0]}
\end{itemize}

\section{软件接口}

\subsection{游戏核心接口}

\texttt{sw/include/fighter\_game.h} 定义了整个游戏状态的数据模型。

核心枚举包括：

\begin{itemize}
  \item \texttt{fighter\_game\_state\_t}: \texttt{MENU / PLAYING / GAME\_OVER}
  \item \texttt{fighter\_visual\_state\_t}: \texttt{IDLE / WALK / JUMP / CROUCH / GUARD / ATTACK / HIT / BLOCK\_STUN / KO}
  \item \texttt{fighter\_attack\_phase\_t}: \texttt{STARTUP / ACTIVE / HIT\_CONFIRM / BLOCK\_CONFIRM / RECOVERY}
  \item \texttt{fighter\_winner\_t} 与 \texttt{fighter\_finish\_reason\_t}
\end{itemize}

核心结构体包括：

\begin{itemize}
  \item \texttt{fighter\_game\_config\_t}：当前默认配置为 \texttt{640x480}，地面 \texttt{y=400}，角色尺寸 \texttt{48x96}，步速 4，起跳速度 -18，重力 1，血量 100，回合时间 \texttt{99*60} 帧。
  \item \texttt{fighter\_player\_state\_t}：保存 \texttt{x/y/vy/hp/facing}、攻击阶段、可视状态、事件标志和状态帧计数。
  \item \texttt{fighter\_game\_t}：保存全局状态、帧计数器、回合计时、胜者、结束原因以及两名玩家状态。
\end{itemize}

核心函数包括：

\begin{itemize}
  \item \texttt{fighter\_game\_init()}
  \item \texttt{fighter\_game\_tick()}
  \item \texttt{fighter\_game\_menu\_animation\_frame()}
  \item \texttt{fighter\_game\_game\_over\_ready()}
  \item \texttt{fighter\_game\_round\_seconds\_remaining()}
\end{itemize}

\subsection{输入接口}

\texttt{sw/include/fighter\_input.h} 定义了输入解析层。

默认键位如下：

\begin{itemize}
  \item \texttt{W}: 跳
  \item \texttt{A}: 左
  \item \texttt{D}: 右
  \item \texttt{S}: 蹲
  \item \texttt{J}: 攻击
  \item \texttt{K}: 防御
  \item \texttt{L}: 退出当前对局
\end{itemize}

组合攻击映射如下：

\begin{itemize}
  \item \texttt{J}: 普通攻击
  \item \texttt{A + J}: \texttt{FIREBALL}
  \item \texttt{D + J}: \texttt{DRAGON\_PUNCH}
  \item \texttt{W + J}: \texttt{JUMP\_ATTACK}
  \item \texttt{W + D + J}: \texttt{FORWARD\_JUMP\_ATTACK}
  \item \texttt{W + A + J}: \texttt{BACK\_JUMP\_ATTACK}
  \item \texttt{S + J}: \texttt{SWEEP}
\end{itemize}

接口分成两层：

\begin{itemize}
  \item \texttt{fighter\_menu\_parser\_*}：面向精细菜单交互，支持 \texttt{START/EXIT} 选项切换。
  \item \texttt{fighter\_player\_parser\_*}：面向战斗输入，输出统一的 \texttt{fighter\_player\_result\_t}。
\end{itemize}

需要注意的是，当前 \texttt{phase1\_demo} 中真正驱动状态机的是 \texttt{fighter\_player\_parser\_*}；菜单解析 API 已经存在，但还没有被完整接入到 \texttt{phase1\_demo} 的主循环。

\subsection{渲染接口}

\texttt{sw/include/fighter\_renderer.h} 定义了当前视频输出抽象。

\begin{itemize}
  \item 优先后端：\texttt{framebuffer}
  \item 回退后端：\texttt{console}
  \item 主入口：\texttt{fighter\_renderer\_init()}、\texttt{fighter\_renderer\_draw()}、\texttt{fighter\_renderer\_close()}
\end{itemize}

渲染器当前会：

\begin{itemize}
  \item 在菜单态加载并缩放两张菜单帧图 \texttt{menu\_frame\_0.ppm} 与 \texttt{menu\_frame\_1.ppm}
  \item 在对局态绘制背景、血条、倒计时、矩形/简化角色图形和调试信息
  \item 在无法打开 \texttt{/dev/fb0} 时退回控制台调试输出
\end{itemize}

\subsection{音频接口}

\texttt{sw/include/fighter\_audio.h} 定义音频命令队列和后端抽象。

音频轨道包括：

\begin{itemize}
  \item \texttt{MENU\_BGM}
  \item \texttt{MENU\_CONFIRM}
  \item \texttt{GAME\_OVER}
\end{itemize}

命令类型包括：

\begin{itemize}
  \item \texttt{PLAY\_ONCE}
  \item \texttt{START\_LOOP}
  \item \texttt{STOP\_LOOP}
\end{itemize}

后端选择逻辑如下：

\begin{itemize}
  \item 默认音频关闭
  \item 只有在 \texttt{phase1\_demo --audio} 时才尝试初始化音频
  \item 初始化顺序为 \texttt{WM8731/MMIO -> command backend fallback}
\end{itemize}

\subsection{USB 键盘采集接口}

\texttt{sw/include/usb\_hid\_keyboard.h} 基于 \texttt{libusb} 实现设备枚举与轮询。

关键点如下：

\begin{itemize}
  \item 最多支持 2 个设备
  \item 只接受 \texttt{HID boot keyboard}
  \item 使用 interrupt IN endpoint 读取 8-byte keyboard report
  \item 设备断开时会自动清零对应 report 并打印错误
\end{itemize}

\section{Verilog 模块接口}

模板中原本希望描述一个“顶层格斗外设 + VGA 子模块”组合，但当前仓库真实存在的项目模块是音频外设，因此本节按当前代码落地情况编写。

\subsection{\texttt{fighter\_audio\_wm8731}}

模块定义见 \texttt{hw/fighter\_audio.sv}。

关键参数：

\begin{itemize}
  \item \texttt{FIFO\_DEPTH=128}
  \item \texttt{CLK\_HZ=50000000}
  \item \texttt{I2C\_RATE\_HZ=100000}
  \item \texttt{XCK\_DIV=4}
  \item \texttt{BCLK\_DIV=16}
  \item \texttt{I2C\_DEVICE\_ADDR=7'h1A}
\end{itemize}

关键端口：

\begin{itemize}
  \item 时钟复位：\texttt{clk}, \texttt{reset\_n}
  \item Avalon-MM：\texttt{avs\_chipselect}, \texttt{avs\_read}, \texttt{avs\_write}, \texttt{avs\_address}, \texttt{avs\_writedata}, \texttt{avs\_readdata}
  \item 音频数据：\texttt{aud\_xck}, \texttt{aud\_bclk}, \texttt{aud\_daclrck}, \texttt{aud\_adclrck}, \texttt{aud\_dacdat}, \texttt{aud\_adcdat}
  \item Codec 配置：\texttt{fpga\_i2c\_sclk}, \texttt{fpga\_i2c\_sdat}
  \item 状态输出：\texttt{codec\_init\_done}, \texttt{codec\_init\_error}
\end{itemize}

\subsection{\texttt{soc\_system\_top}}

板级包装模块见 \texttt{hw/soc\_system\_top.sv}。从项目角度看，它的职责是：

\begin{itemize}
  \item 实例化 \texttt{soc\_system}
  \item 把 HPS USB、DDR3、I2C、UART 等标准管脚接好
  \item 把 \texttt{audio\_*} conduit 连接到 \texttt{AUD\_*}
  \item 暴露 \texttt{audio\_init\_done/error}
  \item 在当前构建里静态关闭未使用 VGA 输出
\end{itemize}

\subsection{当前缺失的模块}

以下模块在设计目标中有明确位置，但当前分支尚未提供实现：

\begin{itemize}
  \item \texttt{vga\_controller}
  \item \texttt{sprite\_engine}
  \item \texttt{ui\_overlay}
  \item \texttt{fighter\_peripheral} 或等价 VGA/MMIO 外设
\end{itemize}

这也是为什么模板里的 VGA 章节必须被改写成“目标架构说明”，不能写成“当前实现”。

\section{实现细节与测试}

\subsection{主循环与运行模式}

\texttt{sw/main\_phase1\_demo.c} 是当前主 demo，支持两种输入模式：

\begin{itemize}
  \item \texttt{--script smoke|ko}
  \item \texttt{--usb}
\end{itemize}

运行时每帧执行：

\begin{enumerate}
  \item 采集或生成输入报告
  \item 解析为 \texttt{fighter\_player\_result\_t}
  \item 执行 \texttt{fighter\_game\_tick()}
  \item 执行 \texttt{fighter\_audio\_process\_commands()}
  \item 执行 \texttt{fighter\_renderer\_draw()}
  \item 用 \texttt{clock\_gettime()} 和 \texttt{nanosleep()} 控制到约 60 FPS
\end{enumerate}

\subsection{已实现的游戏规则}

当前分支已经实现：

\begin{itemize}
  \item 双方左右移动、跳跃、下蹲、防御
  \item 攻击启动、有效、命中确认、格挡确认、恢复
  \item 普通攻击、波、升龙、跳攻、前跳攻、后跳攻、扫腿
  \item 站位朝向自动更新
  \item 角色重叠推开
  \item 命中伤害、格挡削血、受击硬直、格挡硬直
  \item \texttt{KO / DOUBLE\_KO / TIME\_OUT / EXIT}
  \item \texttt{MENU / PLAYING / GAME\_OVER}
\end{itemize}

\subsection{自动测试覆盖}

\texttt{sw/tests/test\_phase1.c} 已覆盖以下关键路径：

\begin{itemize}
  \item 菜单进入与开始对局
  \item 对局中退出返回菜单
  \item KO 进入 \texttt{GAME\_OVER}
  \item 结算动画门控后重新开始
  \item 超时判平与 MMIO 编码
  \item 命中确认状态机
  \item 格挡硬直与 MMIO 字段
  \item 空中越体后的朝向翻转
  \item 地面重叠的分离逻辑
\end{itemize}

\subsection{构建产物}

\texttt{sw/Makefile} 当前会构建：

\begin{itemize}
  \item \texttt{phase1\_demo}
  \item \texttt{phase1\_test}
  \item \texttt{audio\_demo}
  \item \texttt{audio\_probe}
  \item \texttt{input\_demo}，仅在系统安装 \texttt{libusb-1.0} 时构建
\end{itemize}

推荐验证命令：

\begin{lstlisting}[style=fighterstyle]
cd sw
make
./phase1_test
./phase1_demo --script smoke --console
./phase1_demo --script ko --console
./audio_probe
\end{lstlisting}

\subsection{当前风险与限制}

\begin{itemize}
  \item 当前没有 FPGA VGA 格斗渲染模块，模板里的 VGA 硬件框图只能作为后续计划。
  \item \texttt{phase1\_demo} 没有接入精细菜单状态机，当前菜单行为是“任意玩家任意按键开始对局”。
  \item 音频默认关闭，只有显式传入 \texttt{--audio} 才会尝试使用 \texttt{WM8731/MMIO} 或命令行播放器。
  \item 当前没有 Linux kernel driver，显示和音频都走 userspace 直接访问设备节点或物理地址，调试方便，但工程化程度较低。
\end{itemize}

\section{结论与后续工作}

按照当前仓库的真实状态，本项目已经完成了一个可运行的 \texttt{Phase 1} 格斗游戏原型，以及一条可以单独 bring-up 的 \texttt{WM8731} 音频硬件链路。下一阶段最合理的工作不是重写现有 HPS 逻辑，而是把已经固定的 \texttt{fighter\_mmio} 协议接到真正的 FPGA VGA 渲染外设上，这样可以把当前软件侧渲染逐步迁移到课程目标中的 \texttt{HPS logic + FPGA video} 架构。

最直接的增量路线如下：

\begin{enumerate}
  \item 保持 \texttt{fighter\_game.c} 与 \texttt{fighter\_mmio\_encode()} 不变。
  \item 新增一个与 \texttt{fighter\_mmio.h} 完全对齐的 VGA Avalon-MM 从机。
  \item 在 FPGA 侧完成时序、背景、角色和 HUD 的像素输出。
  \item 最后再决定是否补 Linux driver 或继续维持 userspace \texttt{mmap} 方案。
\end{enumerate}

\section{参考文件}

\begin{itemize}
  \item \texttt{sw/main\_phase1\_demo.c}
  \item \texttt{sw/game/fighter\_game.c}
  \item \texttt{sw/render\_if/fighter\_mmio.c}
  \item \texttt{sw/render\_if/fighter\_renderer.c}
  \item \texttt{sw/input/fighter\_input.c}
  \item \texttt{sw/input/usb\_hid\_keyboard.c}
  \item \texttt{sw/audio/fighter\_audio.c}
  \item \texttt{hw/fighter\_audio.sv}
  \item \texttt{hw/fighter\_audio\_hw.tcl}
  \item \texttt{hw/soc\_system\_top.sv}
\end{itemize}

\end{document}
